\scanpage{117}
\mspage{58}
\[
(23)\qquad
\mathcal E_X^{\deleted{\unsure{K}}\inserted{\overline L}}
=\mathcal E_L^{\unsure{g|G}\overline L}
\simeq\mathcal E_U^{\overline L}/\begin{gathered}\text{ss-groupe fermé}\\\text{inv. engendré}\\\text{par les }\unsure{L_{x_i'}|I_{x_i'}}\\(i\in I)\end{gathered}
\]
Notons que $\mathcal E_L^{\overline L}$ est la limite projective des $\mathcal E_U^{\overline L}$, pour des modèles réguliers (affines si on y tient) variables
\[
(24)\qquad\mathcal E_L^{\overline L}\simeq\lim_{\longleftarrow\,U}\mathcal E_U^{\overline L}
\]
et de $=$, bien sûr
\[
(25)\qquad\pi_{L/K}^{\overline L}\simeq\lim_{\longleftarrow\,U}\pi_{U/K}^{\overline L}.
\]
Comment récupérer les $\mathcal E_U^{\overline L}$ \unsure{à} $\mathcal E_L^{\overline L}$ ? Pour \uline{tt} sous-anneau \inserted{d-vr $V$} de \unsure{ess.|en.} de type \unsure{fini} sur $K$ (correspond à un ``germe de modèle'' de $L/K$, au voisinage d'un pt de codim $1$ i.e. du pt gén. d'un diviseur), considé-\newline
rons le point fermé $x$ de $S_x=\operatorname{Spec}(V)$,
\margin{Dans les considérations qui suivent, $K$ (et $X,U$) ne jouent provisoirement aucun rôle --- \unsure{(|[}Mais on pourrait ``plus gén.'' regarder groupes de décomp. \deleted{\unsure{de dimension relative}} d'un des ``$U$'' \uline{mais seulement} \uline{\unsure{épis-comp|épais-comp}} (la ``généralisation'' étant triviale, \unsure{par} $\mathcal E_L\longrightarrow\mathcal E_U$)}
\margin{NB on n'utilise \unsure{les|en} termes ``absolus'' que pour des schémas \uline{épais} (quasi-proj. ?) \uuline{réguliers}}
\margin{116}
